% =====================================================================
\section{Related Work}
\label{sec:related}
% =====================================================================

\paragraph{Defocus deblurring.}
Classical methods estimate a spatially varying defocus map, from edge cues~\cite{zhuo2011defocusmap}
or learned predictors~\cite{lee2019dmenet}, and then deconvolve. Abuolaim and
Brown~\cite{abuolaim2020dpdd} introduced DPDD, real pairs captured with a wide (f/4) and a narrow
(f/22) aperture, together with DPDNet, which uses the dual-pixel (DP) views; later work modelled DP
data further~\cite{abuolaim2021rdpd,xin2021dpdefocus}. Single-image methods design kernel-sharing
atrous convolutions~\cite{son2021kpac}, predict inverse kernels in feature space~\cite{lee2021ifan},
train on light-field-generated defocus~\cite{ruan2022drbnet}, model blur with Gaussian kernel
mixtures or recursive kernels~\cite{quan2021gkmnet,quan2023nrknet}, use large kernels~\cite{ruan2023lakdnet},
or handle the misalignment of aperture-bracketed pairs~\cite{li2023jdrl}; general restoration
backbones~\cite{zamir2022restormer,chen2022nafnet,wang2022uformer} are strong on DPDD as well.
Bokehlicious~\cite{seizinger2025bokehlicious} renders and removes bokeh with aperture control and is
trained on multi-megapixel captures. All of these are trained on crops and evaluated at about $2$--$3$\,MP,
and hence calibrated to the CoC distribution at that resolution. We do not propose another
low-resolution deblurrer; any of them serves as our anchor, and we address the step they leave
open: the native-resolution result.

\paragraph{Ultra-high-definition restoration.}
UHD restoration methods for low-light enhancement, dehazing, deraining and motion
deblurring~\cite{zheng2021uhdmultiguided,wang2023llformer,wang2024uhdformer,zhang2022mcblur}
typically process a strongly downscaled image and recover full resolution with learned upsampling,
bilateral grids or frequency-domain operators. Our decomposition shares this coarse-to-fine spirit,
but it is motivated by the growth of the CoC with resolution, uses an arbitrary existing method as
the coarse stage, and, crucially, exploits the spatially varying amount of detail that survives in
the defocused observation, both locally and in in-focus regions elsewhere, instead of predicting
the fine scale from the coarse estimate alone.

\paragraph{Super-resolution and generative restoration.}
Regression-based super-resolution~\cite{liang2021swinir,chen2023hat} and GAN-based real-world
models~\cite{wang2021realesrgan,zhang2021bsrgan} are natural upsamplers for a low-resolution
deblurring result, and diffusion priors~\cite{saharia2022sr3,whang2022deblurring,xia2023diffir,rombach2022ldm}
adapted from text-to-image models~\cite{wang2024stablesr,lin2024diffbir,yu2024supir,wu2024seesr} have
raised perceptual quality further. Step distillation~\cite{song2023consistency,sauer2024add,yin2024dmd}
has produced one-step restorers~\cite{wang2024sinsr,wu2024osediff}, and diffusion
transformers~\cite{peebles2023dit,esser2024sd3} provide stronger and more scalable priors. These
models are blind to the native observation: they generate detail consistent with the low-resolution
input, which at native resolution is the anchor rather than the photograph. Our experiments show the
consequence, invented text and texture and damaged in-focus regions, in line with the
perception--distortion trade-off~\cite{blau2018perception}. Our diffusion-transformer variant uses
the same one-step recipe but is conditioned on the native observation and restricted to regions
where neither the observation nor the exemplars provide detail.

\paragraph{Reference-based and internal texture transfer.}
Reference-based super-resolution transfers texture from a separate high-resolution reference through
patch correspondence~\cite{zhang2019srntt,yang2020ttsr,jiang2021c2matching,lu2021masasr,cao2022datsr},
and retrieval-augmented restoration conditions large generative models on retrieved
references~\cite{guo2024refir}. Internal learning exploits the recurrence of patches within a single
image, across positions and scales~\cite{glasner2009single,zontak2011internal,shocher2018zssr}, a
principle shared with non-local filtering~\cite{buades2005nlm} and example-based texture
synthesis~\cite{efros1999texture,hertzmann2001analogies,barnes2009patchmatch}. Our exemplar memory
combines both ideas in a setting where neither applies directly: the reference is the observation
itself, its usable parts are defined by the local defocus, and correspondences are established in the
deblurred anchor, where sharp and defocused regions look alike, while texture is taken from the native
observation. We compare against reference-based super-resolution~\cite{cao2022datsr} given the
observation as reference.

\paragraph{Guided upsampling and patch-wise inference.}
Joint bilateral upsampling~\cite{kopf2007jbu}, the guided filter~\cite{he2013guided} and their learned
counterparts~\cite{wu2018deepguided,gharbi2017hdrnet} transfer edges from a high-resolution guide to a
low-resolution solution; they assume a guide that is sharp, whereas ours is sharp only where it is in
focus. Patch-wise inference at high resolution is known to change the behaviour of restoration
networks, partly through global statistics that TLC~\cite{chu2022tlc} corrects at test time, and
tiled generative models fuse overlapping predictions or synchronize them to avoid
seams~\cite{bartal2023multidiffusion,lee2023syncdiffusion,du2024demofusion}. At native resolution we
find seams to be a minor effect compared with the blur-extent gap; the global exemplar memory gives
every tile access to the whole image and makes the same material receive texture from the same
sources across tiles.
